\section{Small positive witnesses}\label{sec:small}

Throughout this section $N\ge N_0=6250$, and
\[
 \rho_0=\frac{30501}{500000},\qquad A_0=\frac{\rho_0}{100}.
\]
We work in the strict-hole configuration of
Lemma~\ref{pr:edge-reduction}, with $0<a\le1/10$ and an actual
simple rightmost zero $b=x+iy$, $x\ge a$, $y\ge0$.
Put $r=|b|$, $p_b=1-r^2$, $\tau=1-y$.
Use $d,S,H,Q$ as in \eqref{pr:field-notation}, and
$D,q,\mu$ as in \eqref{pr:normalized-moments}.

\subsection{The first field bound}

\begin{lemma}\label{sm:normalized-D25}
Let $F(u)=p(i+\rho u)$ be a normalized strict-hole configuration
with an actual simple rightmost zero $\beta=X-i$, where
$\alpha=a/\rho$, $0<\alpha\le X\le1$,
$1/100\le\rho\le9/100$, and $n\ge6251$.
Let $v_j$ list all zeros of $F'$, counted with multiplicity.
Write $h=|F''(\beta)/F'(\beta)|$ and
$d_u=\min_j|\beta-v_j|$. Then $d_u>1$ and $D=d_uh>25$.
\end{lemma}
\begin{proof}
Put
\[
 R=9/100,\qquad Q_u=\sum_j|\beta-v_j|^{-2},\qquad Y=Q_u/h.
\]
The critical domain is
$r_j<\alpha/2$ and
$\rho|v_j|^2/2\le t_j<\rho|v_j-\alpha|^2/2$ for
$v_j=r_j-it_j$. Lemma~\ref{lt:crescent}, scaled back to these
coordinates, gives
\[
 d_u^2>(X-\alpha/2)^2+(1-\eta)^2,\qquad
 \eta=\frac{\rho\alpha^2}{4(1+\sqrt{1-\rho^2\alpha^2/4})}
       \le R\alpha^2/4.
\]
Since $0\le\eta\le R\alpha^2/4<1$, it follows that
\[
 d_u^2>G(\alpha):=(X-\alpha/2)^2+(1-R\alpha^2/4)^2.
\]
This polynomial decreases with $\alpha$ on $0\le\alpha\le X$:
its derivative is
$-X+(1/2-R)\alpha+(R^2/4)\alpha^3<0$.
Hence $d_u^2>1+(1/4-R/2)X^2+(R^2/16)X^4>1$.

Apply the finite polar estimate in Lemma~\ref{or:HB} at
$s=X+1/2$. The inverse Cayley test is strictly to the right of
the original zeros, since $(R(1+1/2)/2)^2<1/2$.
The source root contributes at least
$b_0=4(1-R)/5$. By the full critical comparison
\eqref{or:critical-comparison-formula} and
\eqref{or:small-X-comparison},
\[
 C_{\rm real}<\rho h\mathcal A(X,Y),\qquad
 \mathcal A(X,Y)=
 \begin{cases}1/4,&0<X\le1/4,\\
 X-\kappa(X)Y,&1/4\le X\le1,
 \end{cases}
\]
\[
 \kappa(X)=\frac{X(4X-1)(4X+9)}{40(X+1)}.
\]
Use $f=[1-(9/8)R^2]^{-1}$ and $D_p=1+(3R/4)^2$ in
\eqref{or:parallel-source}. The function $uv/(u+v)$ is increasing
in both variables, so
\begin{equation}\label{sm:normalized-source}
 \rho h\mathcal A(X,Y)>K_0,\qquad
 K_0=\min_{\rho\in\{1/100,R\}}
 \frac1f\left(\frac{6251\rho b_0}{6251\rho+2D_pb_0}-\rho\right)
 =\frac{239012578698741}{379195048000000}>0.
\end{equation}
The expression minimized here is concave in $\rho$
and increases with the degree, as shown by differentiating it in
\eqref{or:direct-K}. Thus the bound holds throughout the scale interval
for every $n\ge6251$. The unphased radial inequality is
$N\rho/2<h(1+Y)$.

If $X\le7/25$, then $\mathcal A\le7/25$, whence
$h>K_0/[R(7/25)]>25$.
If $X\ge7/25$, completion of the square gives
$\mathcal A(X,Y)(1+Y)<31/10$ for every real $Y$.
For an exact proof, put $N_k=X(4X-1)(4X+9)$ and $D_k=40(X+1)$.
The degree-six polynomial
\[
 (62/5)N_kD_k-(XD_k+N_k)^2
\]
has the following positive Bernstein coefficients on $[7/25,1]$:
\[
 \frac{120508031}{244140625},\quad
 \frac{1411226329}{1953125},\quad
 \frac{185744659}{78125},\quad
 \frac{416205433}{78125},\quad
 \frac{30945547}{3125},\quad
 \frac{407837}{25},\quad24527.
\]
Multiplying the two inequalities for the same $X,Y$ therefore gives
$h^2>K_0N/(2(31/10))>625$. Since $d_u>1$, both branches give
$D>25$.
\end{proof}

\begin{lemma}\label{sm:first-D25}
The implication $D>25$ holds on each of the two regions
\[
 \max(x,1-y)>\rho_0,\qquad |b|\le47/50.
\]
Both regions also satisfy $D>70$.
\end{lemma}
\begin{proof}
If $x\ge\rho_0$, Lemma~\ref{lt:direct-D70} gives $D>70$.
We treat $x\le\rho_0$.

Return to physical coordinates. Lemma~\ref{lt:crescent} gives
$d>\tau$. If $1/100\le\tau\le9/100$ and $x\le\tau$, the
preceding lemma applies with $\rho=\tau$, $X=x/\tau$,
$\alpha=a/\tau$; its $D$ is exactly the physical $D$.
For $\tau\ge9/100$, use the matrix envelope in
Proposition~\ref{mf:envelope}. Writing $c=\Re S/H$,
$s=|\Im S|/H$, the rational bound $\pi>333/106$ gives
$\gamma<241/1000$, and its constants satisfy
\[
 Qd^2<21/5+(11/2)Dc+(241/1000)D^2c^2-D^2s^2/2.
\]
Retaining the corresponding negative coefficient $741/1000$ and
completing the square in the radial identity yields
\[
 N<\frac{2Dx}{d}+\frac{p_b}{d^2}
       \left(21/5+(11/2)D+(241/1000)D^2\right)
       +\frac{y^2}{(741/1000)p_b}.
\]
On the whole $x\le\rho_0$, $\tau\ge R=9/100$ range,
\[
 x/d<\rho_0/R,\qquad p_b/d^2<(2-R)/R,\qquad
 y^2/p_b\le(1-R)/(2R).
\]
For the last bound, $x<\tau$ gives $p_b>2\tau(1-\tau)$;
the case $y=0$ is immediate. Under $D\le25$ the preceding upper
bound for $N$ is $640685569/102600<6250$, a contradiction.

Outside the origin box, the remaining small-$x$ range has
$\tau>\rho_0>x$. For $r\le47/50$ one has $\tau\ge.06$.
If $.06\le\tau\le\rho_0$, then
\[
 (.94)^2<\rho_0^2+(1-\rho_0)^2,\qquad
 x^2\le(.94)^2-(1-\tau)^2<\tau^2,
\]
because $\tau^2+(1-\tau)^2$ decreases on this interval.
Thus every case enters either the normalized source or the
large-$\tau$ matrix argument. This proves $D>25$ on both regions.

Proposition~\ref{mf:D25} therefore gives $Qd/H<3$.
Outside the origin box and with $x\le\rho_0$, the radial identity
first gives $D>6250/(2/\rho_0+96)>48$. The matrix envelope excludes
$|s|\ge2/3$, since $c<3/4$, $c^2\le5/9$ would give
\[
 \frac{21/5}{48^2}+\frac{(11/2)(3/4)}{48}
       +\frac{241}{1000}\frac59-\frac29<0.
\]
The phase-sensitive radial bound then gives
\[
 D>\frac{6250\rho_0}{2\rho_0+(4/3)(1-\rho_0)+3(2-\rho_0)}>53.
\]
On $r\le.94$ the initial bound is $D>6250(.06)/8=46.875$.
The matrix envelope excludes $|s|\ge.7$, using $c<.715$ and
$c^2\le.51$. Consequently
\[
 D>\frac{6250(.06)}{2\rho_0+1.4(.94)+3(1.94)}>51.
\]
Both regions satisfy $q<.06$ and $\mu<.02$, so the first
weight-eight row of Proposition~\ref{lt:W8} applies.
Its phase bound $|\Im S|<8Q/H$ and the radial identity give respectively
\[
 D>\frac{6250-48/\rho_0}{98}
   =\frac{27771875}{498183},\qquad
 D>\frac{6250-48/.06}{2\rho_0/.06+97}
   =\frac{27250000}{495167}.
\]
Both exceed $5000/91$, so Proposition~\ref{lt:W8-field} applies.
Its field bound with constant $210$ gives
\[
 D>\frac{6250\rho_0-210}{11/5}=\frac{13701}{176}>70
\]
outside the origin box, and $D>[6250(.06)-210]/(11/5)=75>70$
on $r\le.94$. The geometric estimate covers the large-$x$ branch.
Every numerical comparison increases with $N$, so the full degree
ray is included.
\end{proof}

\subsection{Witnesses below $A_0$}

\begin{proposition}\label{sm:fine}
There is no strict hole with $0<a\le A_0$ and
$\max(x,1-y)>\rho_0$.
\end{proposition}
\begin{proof}
The crescent gives
\[
 d>\rho_0-a/2\ge d_0:=\rho_0-A_0/2.
\]
Lemma~\ref{sm:first-D25} first gives $D>70$.
Then \eqref{lt:field70}, followed in order by
Proposition~\ref{mf:D100}, \eqref{lt:field100}, and
Proposition~\ref{lt:medium139}, gives the strict successive bounds
\[
 D>\frac{6250d_0-596/5}{2(1+A_0)}>100,\qquad
 D>\frac{6250d_0-1899/20}{2(1+A_0)}>140,
\]
\[
 D>\frac{6250d_0-2637/32}{2(1+A_0)}>148.
\]
The last step uses $140>139$.
With $\kappa=11067/5000$, the exact inequality
\[
 (5/4)\kappa^2-(148-23/2)\kappa+296<0,\qquad 2<\kappa<9/4,
\]
and the smaller quadratic branch in Proposition~\ref{lt:medium139}
give $Qd/H<\kappa$.

For $R=td$ define
\[
 \Lambda=\frac{\kappa t}{1-t},\quad
 f=t-\kappa\left(t^2/2+t^3/[3(1-t)]\right),\quad
 B=\frac{\kappa(1-t)}{[1-(1+\kappa)t]^3}.
\]
Thus $H\widehat R=Df$, and the inverse lemma applies with
$\epsilon=(1001/1000)B/148$,
$g=(1001/1000)\nu Bf$, $p_*=(1+\Lambda)^{-1}$ and
$A_*=(1-\Lambda)^{-1}$.
The exact data \cert{FINE6251_INTEGER_DATA.json} give forty-one
triples $(j,m,W)$, with $t=j/200$ and $\nu=m/10000$.
For every row the finite predicates are
\begin{gather}
 0<t,\nu,\Lambda<1,\quad f>0,\quad
 W^2<(1-\nu^2)(148f)^2,\quad\nu(148)f>22/7,\notag\\
 \epsilon(1+e)+ge<1-A_*e,\quad
 \epsilon(1+e)+ge^2<p_*-e,\quad\epsilon+ge<p_*.
 \label{sm:profile-gates}
\end{gather}
Here, for $W=1,\ldots,10$, the rational bounds $e>e^{-W}$ are
\[
 1/2, 1/7, 1/20, 1/50, 1/140, 1/400,
 1/1000, 1/2900, 1/8000, 1/22000.
\]
Their validity follows from the positive exponential sum through
degree one hundred. The data also have increasing $t$ and increasing
$\nu f$. Equations~\eqref{sm:profile-gates} are exactly the full
inverse, corner, four-side and homotopy conditions of
Lemma~\ref{lt:inverse-log}; the resulting cumulative count is
$L_j>\nu_j(7/22)Df_j-2$.

Put $\nu_0f_0=0$ and
\[
 c=\frac{6250}{6251}\left\{
 \frac7{22}\sum_j(1-t_j)^2(\nu_jf_j-\nu_{j-1}f_{j-1})-\frac2{148}
 \right\}>0.
\]
Lemmas~\ref{lt:transport} and \ref{lt:abel}, with $A=A_0$, give
$d<Ja$, where $J=E_{A_0}/[(1-q_{A_0})c]$.
Exact substitution of the forty-one rows gives
$J+1/2<77<100$.
Thus $d<Ja<100a-a/2\le\rho_0-a/2$, contrary to the crescent
lower bound. The finite loss $2/148$ was charged once through
positive cumulative summation. The factor $N/(N+1)$ was bounded
below by $6250/6251$, so the proof holds for every $N\ge6250$.
\end{proof}

\subsection{The low and middle modulus ranges}

\begin{proposition}\label{sm:low-middle}
There is no strict hole with $A_0\le a\le1/10$ and $|b|\le47/50$.
\end{proposition}
\begin{proof}
Lemma~\ref{sm:first-D25} gives $D>70$. Since $d>.06$,
\eqref{lt:field70} gives $D>[6250(.06)-596/5]/(11/5)>100$.
Thus Proposition~\ref{mf:D100} and \eqref{lt:field100} apply.
The simultaneous crescent bound
$d^2>(.06)^2+a^2/4$ gives $D>139$: after positive squaring the
required inequality
\[
 6250\sqrt{(.06)^2+a^2/4}>1899/20+278(1+a)
\]
is exactly
\[
 9688341a^2-(1036801/5)a+613319/400>0.
\]
Its minimum over the entire real line is
$65688734375/155013456>0$.
Consequently Proposition~\ref{lt:medium139} supplies
\begin{equation}\label{sm:middle-input}
 Qd/H<9/4,\qquad |\Im S|<4Q/H,\qquad
 H>\frac{N-C_F/d}{2(1+a)},\qquad C_F=2637/32.
\end{equation}

Put $r_a=|b-a|$. If $r_a=0$, Proposition~\ref{ar:energy}
contradicts the strict hole. Otherwise use the adaptive
shift from Lemma~\ref{bc:blaschke}:
\[
 u=\min(r_a^2,1/2),\quad t=1-u,\quad
 \gamma=b-ta,\quad C=ta[2x-(1+t)a]\ge a^2u(1-u)>0.
\]
It gives
\begin{equation}\label{sm:product}
 -\log B_p(\gamma)+CQ/2<\log(8/C),\qquad
 B_p(\gamma)=\prod_{p(z_i)=0}
              \left|\frac{\gamma-z_i}{1-\overline{z_i}\gamma}\right|.
\end{equation}
A zero product gives a contradiction. Every unselected factor
is at most one.

\smallskip\noindent\emph{The range $r\le.11$.}
Take $d_*=.89$, $V=(9/4)/.89$ and $H_*=6250(.89)/(9/4)$,
using Proposition~\ref{ar:energy} and the prior ratio.
The complete inverse row $R=.1$, $W=8$, $\nu=.99$ satisfies
\eqref{lt:inverse-corner}--\eqref{lt:inverse-sides}. Its count
$K=\nu(7/22)H_*\widehat R-2$ obeys $K/6251>1/100$.
The count lower bound has the form $cN-2$, so its quotient by $n=N+1$
increases. Every counted root has adaptive factor at most
\[
 \frac{R+.1}{1-(.11+R)(.11)}<1/4<e^{-1}.
\]
Hence $-\log B_p>n/100$. Proposition~\ref{ar:variance} gives
\[
 d\le\sqrt{1-(1-.11^2)/n}<1-.49/n.
\]
If $r_a\le.4/n$, the strict hole instead gives
$d>1-r_a\ge1-.4/n$. Thus $r_a>.4/n$, and monotonicity of
$u(1-u)$ on $[0,1/2]$ gives
\[
 C>A_0^2(.4/n)^2[1-(.4/n)^2]>1/(20000000n^2).
\]
The last scalar inequality holds at $n=6251$ and strengthens as $n$
increases. Positive exponential sums give $e^{20}>8(20000000)$
and $e^{10}>6251$. Equation~\eqref{sm:product} would imply
$n/100<20+2\log n$, which fails at $6251$ and thereafter, since
the derivative of the difference is $1/100-2/n>0$.

\smallskip\noindent\emph{The range $.11<r\le.82$.}
Use
\[
 d_*=.18,\quad V=(9/4)/.18,\quad
 H_*=[6250-C_F/.18]/2.2,\quad R=11/250,\quad\nu=.99,\quad W=6.
\]
The full inverse predicates hold. If $z=b+v$ is counted,
$\Re v\le0$ by rightmostness. Therefore
$|\gamma-z|^2=|v+ta|^2\le R^2+.01$ and the denominator is
at least $1-(.82+R)(.82)>0$. The exact squared comparison bounds
the factor by $3/8$. Since $e^{.98}<8/3$, the count gives
$-\log B_p>.98K>21.5$.
Also $r_a>.11-a$ and
\[
 C>a^2(.11-a)^2[1-(.11-a)^2]>4\cdot10^{-9}.
\]
Indeed the concave function $a(.11-a)$ exceeds $1/15000$ at
both endpoints of $[A_0,.1]$, and the remaining factor exceeds
$.9$. The exact positive sum $e^{21.5}>8/(4\cdot10^{-9})$
contradicts \eqref{sm:product}.

\smallskip\noindent\emph{The range $.82\le r\le.94$.}
Here $r_a^2\ge(.82-.1)^2>1/2$, so the shift is
$\gamma=b-a/2$ and $C=ax-3a^2/4\ge a^2/4$.
We use twenty-three product rows, specified by rational inputs in
\cert{LOW_MIDDLE6251_PRODUCT_PROPOSAL.json}.
For each row $B_L\le r\le B_U$, $a_L\le a\le a_U$, $x\le X$,
choose the displayed $d_*,V,H_*$ from that data, and nested
physical radii $R_j$. Each full contour uses
\[
 \Lambda_j=\frac{VR_j}{1-R_j/d_*},\quad
 \widehat R_j=R_j-V\left(\frac{R_j^2}{2}
                         +\frac{R_j^3}{3(d_*-R_j)}\right),\quad
 B_j=\frac{V}{(1-R_j/d_*)^2(1-\Lambda_j)^3},
\]
with $\epsilon_j=(1001/1000)B_j/H_*$ and
$g_j=(1001/1000)\nu_jB_j\widehat R_j$.
All predicates \eqref{lt:inverse-corner}--\eqref{lt:inverse-sides}
are checked at $H_*$, together with
$\nu_jH_*\widehat R_j>22/7$.
Lemma~\ref{lt:half-disk-factor} supplies a whole-disk factor $f_j<1$.
The weights $w_j<-\log f_j$ are rational lower bounds from the
positive atanh series after dyadic scaling into $[1/2,1)$.
They decrease with radius. All count slopes increase, and positive
Abel summation \eqref{lt:abel-product} gives the budget
\[
 \mathcal B=H_*\sum_jw_j(s_j-s_{j-1})-2w_1,\qquad
 s_j=\nu_j(7/22)\widehat R_j,\quad s_0=0.
\]
For a positive whole-row lower bound $C_{\min}$ and a rational
target $T_*$, the final predicates are
\begin{equation}\label{sm:middle-budget}
 e^{T_*}>8/C_{\min},\qquad
 \mathcal B+6250C_{\min}/2>T_*.
\end{equation}
The energy term follows from Proposition~\ref{ar:energy} and
$C\ge C_{\min}$. Exponential bounds are finite positive sums or
\eqref{lt:exp-upper}; radical enclosures are accepted by exact squares.

For $.82\le r\le.925$ take $d_*=.075$, $V=(9/4)/d_*$,
$H_*<[6250-C_F/d_*]/[2(1+a_U)]$, $X=.925$,
$C_{\min}=a_L^2/4$. The seven closed witness intervals have endpoints
\[
 A_0, .003, .01, .02, .04, .06, .08, .1.
\]
The respective strict reserves in \eqref{sm:middle-budget} exceed
$3.03,4.54,4.41,1.91,1.45,1.79,2.99$.

For $.925\le r\le.94$, $x\ge\rho_0$, take $X=.94$ and choose
$d_*$ below at least one of
\[
 .06,\quad \rho_0-a_U/2,\quad \sqrt{(.06)^2+a_L^2/4}.
\]
Each displayed quantity is a crescent lower bound for $d$. Choose
$V\ge\min((9/4)/d_*,1/(\rho_0-a_U/2))$ and
$H_*<[6250-C_F/d_*]/[2(1+a_U)]$.
The lower bound $C_{\min}$ is the exact minimum of
$a\max(a,\rho_0)-3a^2/4$ on the full row. Below $\rho_0$
the function is concave; above it the function is $a^2/4$ and
increases. The eleven intervals have endpoints
\[
 A_0, .003, .01, .02, .03, .04, .05, .06, .07, .08, .09, .1.
\]
All eleven strict reserves exceed $.067$.

It remains to consider $.925\le r\le.94$, $x\le\rho_0$.
Put $Y=\tau Q/H$. If $Y\le5/4$, the phase and radial identity give
\begin{equation}\label{sm:low-moment-field}
 H>\frac{N-8(5/4)/.06}{2\rho_0+(2-.06)(5/4)}.
\end{equation}
This uses $\tau\ge.06$, $p_b/\tau\le2-\tau\le1.94$, and
$|\Im S|<4Q/H$.
Take $d_*=.06$, $V=(5/4)/.06$, an integer $H_*$ below
\eqref{sm:low-moment-field}, $X=\rho_0$ and $C_{\min}=a_L^2/4$.
The five intervals with endpoints
$A_0,.003,.01,.02,.04,\rho_0$ have reserves exceeding
$3.68,5.43,5.44,3.01,2.31$.
These twenty-three rows comprise 1807 nested complete contours;
together with the two earlier single contours, there are 1809.

Finally, in this small-$x$ region $\tau<.078$ since
$y^2\ge.925^2-\rho_0^2>(1-.078)^2$.
If $\tau\ge\rho_0$ then $x\le\tau$; if $.06\le\tau\le\rho_0$
the seam used in Lemma~\ref{sm:first-D25} again gives $x<\tau$.
The normalization $u=(z-i)/\tau$ therefore gives $T=1$,
$0<\alpha=a/\tau\le X=x/\tau\le1$, $d_u>1$ and
$Y=Q_u/h<9/4$. The only branch not covered by the preceding
product rows lies in
\[
 .06\le\tau\le.078,\qquad 0<\alpha\le X\le1,\qquad
 5/4\le Y\le9/4.
\]
The bounds $D>139$ and \eqref{sm:middle-input} supply the hypotheses
of Lemma~\ref{sm:T-moment-extension}. In normalized coordinates the
field bound reads $h>[N\tau-C_F]/[2(1+\tau)]$, since $d_u>1$
and $a\le\tau$. The lemma excludes this entire normalized range.

All modulus and witness seams are closed. Each field lower bound
is affine increasing in $N$, while every fixed count and energy
lower bound strengthens along the degree ray. The only logarithmic
degree budget, in the low range, was checked with its derivative.
Thus every configuration in the proposition is covered.
\end{proof}

\subsection{The high modulus range}

\begin{proposition}\label{sm:high}
There is no strict hole with $A_0\le a\le1/10$, $|b|>.94$,
outside the origin box $\max(x,1-y)\le\rho_0$.
\end{proposition}
\begin{proof}
If $x<\rho_0$, the exact comparisons
\[
 .94^2>(.0433)^2+(1-\rho_0)^2,\qquad
 .94^2-\rho_0^2>(1-.062)^2,\qquad .0433/.062>.69
\]
give $\rho_0<\tau<.062$ and $.69<x/\tau<1$.
The normalization $u=(z-i)/\tau$ therefore lies entirely in
Lemma~\ref{sm:T-highX-extension}, with $T=1$,
$\rho_0\le\tau\le.062$, $.69\le X\le1$.
We may henceforth assume $x\ge\rho_0$.

Lemma~\ref{lt:direct-D70} gives $D>70$ geometrically.
Proposition~\ref{lt:D70} then gives phase $K_0=21/4$ and the
field \eqref{lt:field70}. Since $|b-a|>.84$, the adaptive shift
is $\gamma=b-a/2$, and \eqref{sm:product} holds with
$C=ax-3a^2/4>0$.
Put $q_N=Q/N$. Cauchy--Schwarz and the radial identity give
\begin{equation}\label{sm:phase-product}
 1<2x\sqrt{q_N}+p_bq_N+2K_0Q/(NH).
\end{equation}

\smallskip\noindent\emph{The first witness cap.}
Suppose $a\ge a_0$ and $x\ge x_L>3a_0$. Concavity of
$a\mapsto ax-3a^2/4$ on $[a_0,x]$ gives
$C\ge a_0(x-3a_0/4)$: the other endpoint is $x^2/4$, which
is larger when $x>3a_0$.
Choose $L$ with $e^L>8/[a_0(x_L-3a_0/4)]$, and put
$s=3a_0/4$, $k=2L/(N_0a_0)$.
The product inequality gives $q_N<k/(x-s)$.
The reciprocal half-plane gives $Q/H<2/x$.
Thus \eqref{sm:phase-product} requires
\begin{equation}\label{sm:initial-cap-function}
 1<2x\sqrt{\frac{k}{x-s}}+
       (1-B_L^2)\frac{k}{x-s}+\frac{4K_0}{N_0x}
       \qquad(r\ge B_L).
\end{equation}
This expression is convex as a function of $v=\sqrt{x-s}$:
its terms are $2\sqrt k(v+s/v)$,
$(1-B_L^2)k/v^2$, and $4K_0/[N_0(v^2+s)]$.
The final term has nonnegative second derivative when $3v^2\ge s$.
It suffices to check the two endpoints of the full interval in $x$.
For
\[
 B_L=.94,\quad x_L=\rho_0,\quad x_U=1,\quad
 a_0=.013,\quad L=1879/200,\quad K_0=21/4,
\]
both endpoint values are strictly below one. These are rational
checks: the remainder after subtracting the last two terms is
positive, and its square exceeds $4x^2k/(x-s)$.
The logarithmic comparison uses a positive exponential sum.
We conclude $a<.013$.

The cap $a<.013$ gives $d>\rho_0-.0065=.054502$, so
\eqref{lt:field70} gives $D>100$.
Thus Proposition~\ref{mf:D100} and \eqref{lt:field100} apply.
For $r\ge.946$,
Lemma~\ref{lt:same-gap-field}, with $a_U=.013$, $x_0=\rho_0$,
$K_0=21/4$ and $p_U=1-.946^2$, gives
\[
 D>321306321875/1995240393>160.
\]
Proposition~\ref{lt:medium139} therefore supplies phase four,
field constant $C_F=2637/32$, and the smaller ratio branch $Y(D)$
throughout that modulus range.

\smallskip\noindent\emph{Eight successive Schur estimates.}
For a known witness cap $a\le a_U$ and modulus floor $B_L$,
Lemma~\ref{lt:same-gap-field} with phase four gives a whole-region
lower bound for $D$. Choose $D_0$ strictly below that value, then
choose $\kappa$ with
\[
 2<\kappa<9/4,\qquad
 (5/4)\kappa^2-(D_0-23/2)\kappa+2D_0<0.
\]
The smaller branch implies $Qd/H<\kappa$.
Use the complete Schur profiles as in
\eqref{sm:profile-gates}, replacing $148$ and $11067/5000$ by
$D_0,\kappa$. Put $\nu_0f_0=0$. If
\[
 c=\frac{6250}{6251}\left\{
 \frac7{22}\sum_j(1-t_j)^2(\nu_jf_j-\nu_{j-1}f_{j-1})
 -\frac2{D_0}\right\}>0,
\]
then, with $J=E_{a_U}/[(1-q_{a_U})c]$, Lemma~\ref{lt:abel}
gives $d<Ja$ and hence
$x<(J+1/2)a$. The complete sequence is
\begin{center}\small
\begin{tabular}{rrrrrrr}
\toprule
$B_L$&prior $a_U$&$D_0$&$\kappa$&$M$&upper bound for $J+1/2$&rows\\\midrule
$.946$&$.013$&$162$&$2.1928$&$75$&$74.260$&$42$\\
$.955$&$.01108$&$198$&$2.1545$&$67$&$66.009$&$44$\\
$.97$&$.00779$&$301$&$2.0985$&$57$&$56.441$&$49$\\
$.98$&$.00572$&$442$&$2.0659$&$52$&$51.715$&$52$\\
$.99$&$.00430$&$781$&$2.0367$&$49$&$48.244$&$54$\\
$.995$&$.00296$&$1256$&$2.0226$&$47$&$46.493$&$56$\\
$.997$&$.00216$&$1651$&$2.0172$&$46$&$45.798$&$57$\\
$.998$&$.00173$&$1951$&$2.0145$&$46$&$45.515$&$57$\\\bottomrule
\end{tabular}
\end{center}
All 411 triples $(j,m,W)$ are rational integer data in
\cert{HIGH6251_INTEGER_DATA.json}, with $t=j/200$,
$\nu=m/10000$. For each, the complete inverse, every side and
corner, the amplitude homotopy, positive odd height, and positive
Abel increment are verified by \eqref{sm:profile-gates}.
Thus the conclusion of each row is $x<Ma$ on its whole domain.

After each row we improve the witness cap, in the same order.
Suppose $a\ge a_0$ and $x<Ma$, and put $x_0=Ma_0$,
$\gamma_M=1/M-3/(4M^2)$. Concavity gives
\begin{equation}\label{sm:joint-C}
 C\ge
 \begin{cases}
 a_0(x-3a_0/4),&\rho_0\le x\le x_0,\\
 \gamma_Mx^2,&x\ge x_0.
 \end{cases}
\end{equation}
The two expressions agree at $x_0$. For the first branch use
$k=2L/(N_0a_0)$, $s=3a_0/4$, and for the second use
$k_q=2L/(N_0\gamma_M)$. The cap $a<a_U$ gives
the phase error $8/[N_0(x-a_U/2)]$ in
\eqref{sm:phase-product}. The first branch is convex in
$\sqrt{x-s}$ when $3(x-s)>s-a_U/2$; the second branch requires
\[
 1<2\sqrt{k_q}+(1-B_L^2)k_q/x^2+\frac8{N_0(x-a_U/2)},
\]
whose right side decreases with $x$.
Hence positive squared endpoint checks at $x=\rho_0,x_0$, together
with $e^L>8/[a_0(\rho_0-3a_0/4)]$, exclude $a\ge a_0$.
The rows are
\begin{center}\small
\begin{tabular}{rrrr}
\toprule
$B_L$&$M$&new cap $a_0$&$L$\\\midrule
$.946$&$75$&$.01108$&$9.528$\\
$.955$&$67$&$.00779$&$9.834$\\
$.97$&$57$&$.00572$&$10.115$\\
$.98$&$52$&$.00430$&$10.382$\\
$.99$&$49$&$.00296$&$10.738$\\
$.995$&$47$&$.00216$&$11.043$\\
$.997$&$46$&$.00173$&$11.26$\\
$.998$&$46$&$.00146$&$11.426$\\\bottomrule
\end{tabular}
\end{center}
The prior cap of each Schur row is exactly the conclusion of the
preceding cap row, and its modulus domain is contained in the
preceding one. Thus these estimates form a finite ordered chain.

\smallskip\noindent\emph{Root counts in physical disks.}
The closed initial domains are as follows. Each contains
all $x$ from $\rho_0$ to the stated upper endpoint, and all
$a\in[A_0,a_U]$ compatible with the configuration.
\begin{center}\small
\begin{tabular}{crrrr}
\toprule
modulus interval&$a_U$&$M$&upper $x$&closed leaves\\\midrule
$[.94,.946]$&$.013$&--&$1$&1\\
$[.946,.955]$&$.01108$&75&$.831$&1\\
$[.955,.97]$&$.00779$&67&$.52193$&1\\
$[.97,.98]$&$.00572$&57&$.32604$&1\\
$[.98,.99]$&$.00430$&52&$.2236$&2\\
$[.99,.995]$&$.00296$&49&$.14504$&3\\
$[.995,.997]$&$.00216$&47&$.10152$&4\\
$[.997,.998]$&$.00173$&46&$.07958$&4\\
$[.998,.9985]$&$.00146$&46&$.06716$&3\\\bottomrule
\end{tabular}
\end{center}
The data give eleven midpoint splits, each along the modulus or
$x$ coordinate, with both closed children retained. This gives
twenty closed leaves. The binary rule regenerates the whole cover,
including the boundaries of all initial domains and their common seams.

On a leaf $B_L\le r\le B_U$, $x_L\le x\le x_U$, put
$d_*=x_L-a_U/2$, $V=1/d_*$ and $p_{\max}=1-B_L^2$.
The available field bounds at the degree floor are
\begin{align}
 H_{\mathrm{hole}}&=\frac{6250-C_{\mathrm{field}}/d_*}{2(1+a_U)},\notag\\
 H_{\mathrm{phase}}&=\frac{6250-2K_0/d_*}
 {\max\{2x_L+p_{\max}/(x_L-a_U/2),
          2x_U+p_{\max}/(x_U-a_U/2)\}}.
 \label{sm:high-leaf-field}
\end{align}
The denominator in the second bound is the maximum of a convex
function on the full interval. In the first initial domain use
$K_0=21/4$, $C_{\mathrm{field}}=1899/20$; in all others use
$K_0=4$, $C_{\mathrm{field}}=2637/32$.
All numerators are positive. Each leaf declares a field mode and an
integer $H_*$ strictly below the corresponding expression in
\eqref{sm:high-leaf-field}.
Its full physical radius $R$, $\nu=m/10000$ and integer $W$
satisfy the inverse predicates with $d_*,V,H_*$ as above.
Thus the actual cumulative count is greater than
$\nu(7/22)H_*\widehat R-2$.

For its factor $f<1$ the leaf verifies
\eqref{lt:half-disk-gate} with $P=1-B_U^2$,
$T=a_U/(2P)$, $U=R/P$ and $L=1+2Tx_U$.
In the first initial domain take $C_{\min}=A_0(x_L-3A_0/4)$.
In every other initial domain use both lower constraints on
$a$:
\[
 C_{\min}=\max\{A_0(x_L-3A_0/4),\gamma_Mx_L^2\}.
\]
Every field in \eqref{sm:high-leaf-field} has the form $(N-c)/F$
with $c\ge0$. Equation~\eqref{lt:energy-ray} therefore gives
$Q>H_*^2/6250$ throughout the full degree ray. For a rational
$\ell<-\log f$ the complete product lower budget is
\begin{equation}\label{sm:high-budget}
 \ell\{\nu(7/22)H_*\widehat R-2\}
       +\frac{C_{\min}H_*^2}{12500}.
\end{equation}
Each leaf has a rational target $T_*$ with
$e^{T_*}>8/C_{\min}$ and with
\eqref{sm:high-budget} strictly greater than $T_*$.
The logarithmic lower weight is checked by
$\mathcal E_{24}(\ell)<1/f$, and the target by the positive
exponential sum through degree one hundred. All twenty whole-leaf
reserves exceed $.2558$. The factor estimate holds for every
original zero in each physical disk, and all remaining product
factors are at most one.

\smallskip\noindent\emph{The unit-circle endpoint.}
On $r\ge.9985$, the preceding chain gives $x<46a$ and
$a<.00146$. Hence $C\ge\gamma_{46}x^2$.
Take $L=11521/1000$; the exact positive exponential gate gives
$e^L>8/(\gamma_{46}\rho_0^2)$.
The product, phase and radial identities would require
\[
 1<2\sqrt{k_q}+(1-.9985^2)k_q/x^2
                 +\frac8{6250(x-.00146/2)},\qquad
 k_q=\frac{2L}{6250\gamma_{46}}.
\]
The right side decreases for all $x\ge\rho_0$. At $x=\rho_0$
it is less than $.990541<1$; equivalently the positive remainder
after the two rational terms has square exceeding $4k_q$ by more
than $.0158$. This excludes the whole endpoint range, including
$p_b=0$ and $|b|=1$.

The fixed input file contains all eight Schur profiles, the eight
successive caps, the nine initial domains, all eleven exact splits,
the twenty leaf witnesses and the endpoint parameters.
\cert{check_high6251_exact.py} reconstructs all the displayed
rational predicates and every inherited-domain equality.
The first fields and the subsequent same-gap bounds increase with
$N$; the Schur factor uses the smallest $N/(N+1)$; and the counts
and field-energy budgets strengthen by \eqref{lt:energy-ray}.
Thus the proof covers every $N\ge6250$.
\end{proof}
